\section{Conclusion}

Matched-Budget Component Attribution asks where, specifically, a learned component adds value inside a
rule-based system, holding everything but the insertion point fixed. Applied to one strategy in three
markets (Study 1) and then, pre-registered, to five historically documented strategies on
\STwoInstruments{} instruments across \STwoMarkets{} markets and two independent freezes (Study 2), it
gives a consistent answer: at a matched budget of generic price features, machine learning rarely
changes a rule-based strategy significantly (\STwoMktSigGain{} significant gains and \STwoMktSigLoss{}
losses in \STwoMktTests{} tests); narrowing components are small, mostly positive, transferable, and
safe; replacing components are high-variance and help only when the replaced rule has a structural
weakness, as ML exit does for RSI(2) mean reversion; and realistic execution matters more than any
component. Because every number in this paper is generated from one pre-registered configuration by the
released pipeline, each of these claims can be re-derived --- or overturned --- by re-running it.
